%%%%%%%%%%%%%%%%%%% use for abstract
%%% refer q31.tex
\NoBgThispage
\begin{abstract}

The Smart Queue Management System is a web-based application developed to digitize and automate the process of queue management in service-based organizations such as hospitals, banks, and government offices. Traditional queue management methods rely on physical token machines and manual processes, which lead to long waiting times, poor customer experience, and inefficient service delivery. This project presents a practical digital solution to replace the conventional approach.

The system is developed using a three-tier client-server architecture. ReactJS is used for the frontend user interface, ASP.NET Core Web API (.NET 10) is used for the backend REST API, and PostgreSQL is used as the relational database. The system allows customers to generate digital queue tokens by entering their name, view the currently serving token, and monitor the real-time queue status. The queue status page auto-refreshes every five seconds to provide near-real-time updates. An admin module enables service staff to log in, call the next token, complete tokens, and reset the queue when required.

Entity Framework Core is used as the Object-Relational Mapper for database access, which enables automatic creation of database tables on the first run without any manual SQL setup. Cross-Origin Resource Sharing (CORS) is enabled on the backend to allow communication between the frontend and backend running on different ports. The system is designed to operate entirely on a local network, making it cost-effective and independent of internet connectivity.

Testing confirmed that all key features function correctly, including sequential token generation, real-time status updates, admin authentication, and complete queue lifecycle management. The system is simple, beginner-friendly, and serves as a practical demonstration of modern full-stack web development.

{\bf Keywords}: {\it Queue Management, Token Generation, ReactJS, ASP.NET Core, PostgreSQL, REST API, Web Application, MCA Project}

\end{abstract}

